\ifx\allfiles\undefined
\documentclass[12pt, a4paper, oneside, UTF8]{ctexbook}
\def\path{../config}
\input{\path/package.tex}

% 定理环境设置
\input{\path/theorem1_zh.tex}

\input{\path/custom.tex}

% 修改参数改变封面样式，0 默认原始封面、内置其他1、2、3种封面样式
\def\myIndex{3}


\ifnum\myIndex>0
    \input{\path/cover_package_\myIndex} 
\fi

\def\myTitle{线性代数笔记}
\def\myAuthor{M.K.}
\def\myDateCover{\today}
\def\myDateForeword{\today}
\def\myForeword{前言}
\def\myForewordText{
    
    这是一个基于\LaTeX{}模板的线代笔记。

}
\def\mySubheading{副标题}


\begin{document}
% \input{\path/cover_text_\myIndex.tex}

\newpage
\thispagestyle{empty}
\begin{center}
    \Huge\textbf{\myForeword}
\end{center}
\myForewordText
\begin{flushright}
    \begin{tabular}{c}
        \myDateForeword
    \end{tabular}
\end{flushright}

\newpage
\pagestyle{plain}
\setcounter{page}{1}
\pagenumbering{Roman}
\pagestyle{fancy}
\fancyfoot[C]{\thepage}
\renewcommand{\headrulewidth}{0.4pt}
\renewcommand{\footrulewidth}{0pt}
\tableofcontents

\newpage
\pagenumbering{arabic}
\setcounter{page}{1}








\else
\fi
\chapter{向量空间与线性变换}
\section{线性相关与线性无关}
\begin{definition}
    设$\bm{V}$是域$\bm{F}$上的向量空间，向量组$\bm{\beta},\bm{\alpha}_1, \bm{\alpha}_2, \cdots, \bm{\alpha}_n \in \bm{V}$，若\textcolor{Burgundy}{存在}数$k_1, k_2, \cdots, k_n \in \bm{F}$，使得
    \[
        k_1\bm{\alpha}_1 + k_2\bm{\alpha}_2 + \cdots + k_n\bm{\alpha}_n = \bm{\beta}
    \]
    则称向量$\bm{\beta}$是$\bm{\alpha}_1, \bm{\alpha}_2, \cdots, \bm{\alpha}_n$的\textbf{线性组合}，
    或称$\bm{\beta}$可由向量组$\bm{\alpha}_1, \bm{\alpha}_2, \cdots, \bm{\alpha}_n$\textbf{线性表示}，
    此时称$k_1, k_2, \cdots, k_n$为线性系数或表示系数．
\end{definition}
\begin{corollary}
    设$\bm{V}$是域$\bm{F}$上的向量空间，向量组$\bm{\alpha}_1, \bm{\alpha}_2, \cdots, \bm{\alpha}_n \in \bm{V}$，则$\bm{0}$一定能由向量组$\bm{\alpha}_1, \bm{\alpha}_2, \cdots, \bm{\alpha}_n$线性表示，即
    \[
        0\bm{\alpha}_1 + 0\bm{\alpha}_2 + \cdots + 0\bm{\alpha}_n = \bm{0}
    \]
\end{corollary}
\begin{definition}
    设$\bm{V}$是域$\bm{F}$上的向量空间，向量组$\bm{\alpha}_1, \bm{\alpha}_2, \cdots, \bm{\alpha}_n \in \bm{V}$，若存在\textcolor{Burgundy}{不全为零}的数$k_1, k_2, \cdots, k_n \in \bm{F}$，使得
    \[
        k_1\bm{\alpha}_1 + k_2\bm{\alpha}_2 + \cdots + k_n\bm{\alpha}_n = \bm{0}
    \]
    则称向量组$\bm{\alpha}_1, \bm{\alpha}_2, \cdots, \bm{\alpha}_n$\textbf{线性相关}，否则称\textbf{线性无关}．
\end{definition}
\begin{theorem}
    设$\bm{V}$是域$\bm{F}$上的向量空间，向量组$\bm{\alpha}_1, \bm{\alpha}_2, \cdots, \bm{\alpha}_n \in \bm{V}$，则向量组$\bm{\alpha}_1, \bm{\alpha}_2, \cdots, \bm{\alpha}_n$线性相关的充要条件是其中至少有一个向量能由其余向量线性表示．
\end{theorem}
\begin{theorem}
    设有$n\times m$矩阵$\bm{A}$，其列向量组线性相关的充要条件是$\rank(\bm{A}) < m$．
\end{theorem}
\begin{theorem}
    设$\bm{V}$是域$\bm{F}$上的向量空间，向量组$\bm{\alpha}_1, \bm{\alpha}_2, \cdots, \bm{\alpha}_n \in \bm{V}$，若$\operatorname{dim}\mathrm{Span}\{\bm{\alpha}_1, \bm{\alpha}_2, \cdots, \bm{\alpha}_n\}<n$，则向量组$\bm{\alpha}_1, \bm{\alpha}_2, \cdots, \bm{\alpha}_n$线性相关．
\end{theorem}
\begin{theorem}
    线性齐次方程组$\bm{A}\bm{x}=\bm{0}$有非零解的充要条件是$\bm{A}$的列向量线性相关．
\end{theorem}
\section{向量组}
\begin{definition}
    设$\bm{V}$是域$\bm{F}$上的向量空间，向量组$\bm{\alpha}_1, \bm{\alpha}_2, \cdots, \bm{\alpha}_n \in \bm{V}$，则由向量组$\bm{\alpha}_1, \bm{\alpha}_2, \cdots, \bm{\alpha}_n$的全体线性组合所构成的集合，称为由向量组$\bm{\alpha}_1, \bm{\alpha}_2, \cdots, \bm{\alpha}_n$所生成的子空间，记为
    \[
        \mathrm{Span}\{\bm{\alpha}_1, \bm{\alpha}_2, \cdots, \bm{\alpha}_n\}
    \]
\end{definition}
\begin{definition}
    设$\bm{V}$是域$\bm{F}$上的向量空间，向量组$\bm{\alpha}_1, \bm{\alpha}_2, \cdots, \bm{\alpha}_n \in \bm{V}$，若向量组$\bm{\alpha}_1, \bm{\alpha}_2, \cdots, \bm{\alpha}_n$线性无关，且任一向量$\bm{\beta} \in \bm{V}$均能由向量组$\bm{\alpha}_1, \bm{\alpha}_2, \cdots, \bm{\alpha}_n$线性表示，则称向量组$\bm{\alpha}_1, \bm{\alpha}_2, \cdots, \bm{\alpha}_n$为$\bm{V}$的一个极大线性无关组．
\end{definition}
\begin{definition}
    设$\bm{V}$是域$\bm{F}$上的向量空间，向量组$\bm{\alpha}_1, \bm{\alpha}_2, \cdots, \bm{\alpha}_n \in \bm{V}$，则向量组$\bm{\alpha}_1, \bm{\alpha}_2, \cdots, \bm{\alpha}_n$的极大线性无关组的向量个数，称为向量组$\bm{\alpha}_1, \bm{\alpha}_2, \cdots, \bm{\alpha}_n$的秩，记为$\rank(\bm{\alpha}_1, \bm{\alpha}_2, \cdots, \bm{\alpha}_n)$．
\end{definition}
\begin{theorem}
    设$\bm{A}$为$m\times n$矩阵，则$\bm{A}$的行秩等于列秩，即
    \[
        \rank(\bm{A}) = \rank(\bm{A}^T)
    \]
\end{theorem}
\begin{myRemark}
    向量组的秩相等，向量组不一定等价，线性方程组不一定同解
\end{myRemark}
\section{向量空间}
\begin{definition}[向量空间]
设 $V$ 是一个非空集合，$F$ 是一个域．如果在 $V$ 上定义了向量加法和标量乘法，且满足以下公理，则称 $V$ 为 $F$ 上的\textbf{向量空间}：
\begin{enumerate}
    \item $\bm{u} + \bm{v} = \bm{v} + \bm{u}$ (加法交换律)
    \item $(\bm{u} + \bm{v}) + \bm{w} = \bm{u} + (\bm{v} + \bm{w})$ (加法结合律)
    \item $\exists \bm{0} \in V, \forall \bm{u} \in V, \bm{u} + \bm{0} = \bm{u}$ (零向量)
    \item $\forall \bm{u} \in V, \exists -\bm{u} \in V, \bm{u} + (-\bm{u}) = \bm{0}$ (负向量)
    \item $a(b\bm{u}) = (ab)\bm{u}$
    \item $k(\bm{u} + \bm{v}) = k\bm{u} + k\bm{v}$
    \item $(a+b)\bm{u} = a\bm{u} + b\bm{u}$
    \item $1\bm{u} = \bm{u}$
\end{enumerate}
\end{definition}
\begin{definition}
    设$\bm{V}$是域$\bm{F}$上的向量空间，向量组$\bm{\alpha}_1, \bm{\alpha}_2, \cdots, \bm{\alpha}_n \in \bm{V}$，若向量组$\bm{\alpha}_1, \bm{\alpha}_2, \cdots, \bm{\alpha}_n$线性无关，且其生成的子空间$\mathrm{Span}\{\bm{\alpha}_1, \bm{\alpha}_2, \cdots, \bm{\alpha}_n\}=\bm{V}$，则称向量组$\bm{\alpha}_1, \bm{\alpha}_2, \cdots, \bm{\alpha}_n$为向量空间$\bm{V}$的一组基底，简称\textbf{基}．
\end{definition}
\begin{definition}
    设$\bm{V}$是域$\bm{F}$上的向量空间，若$\bm{V}$有一组有限基，则称$\bm{V}$为有限维向量空间，且称该组基的向量个数为$\bm{V}$的\textbf{维数}，记为$\operatorname{dim}\bm{V}$；否则称$\bm{V}$为无限维向量空间．
\end{definition}
\begin{definition}
    设$\bm{V}$是域$\bm{F}$上的$n$维向量空间，$\bm{e}_1, \bm{e}_2, \cdots, \bm{e}_n$是$\bm{V}$的一组基，向量$\bm{\alpha} \in \bm{V}$，则存在唯一的一组数$k_1, k_2, \cdots, k_n \in \bm{F}$，使得
    \[
        \bm{\alpha} = k_1\bm{e}_1 + k_2\bm{e}_2 + \cdots + k_n\bm{e}_n
    \]
    称数$k_1, k_2, \cdots, k_n$为向量$\bm{\alpha}$在基$\bm{e}_1, \bm{e}_2, \cdots, \bm{e}_n$下的坐标．
\end{definition}
\begin{theorem}
    设$\bm{V}$是域$\bm{F}$上的$n$维向量空间，$\bm{e}_1, \bm{e}_2, \cdots, \bm{e}_n$和$\bm{e}'_1, \bm{e}'_2, \cdots, \bm{e}'_n$是$\bm{V}$的两组基，
    则存在可逆矩阵$\bm{P}$，使得
    \[
        \begin{bNiceMatrix}
            \bm{e}'_1 & \bm{e}'_2 & \cdots & \bm{e}'_n
        \end{bNiceMatrix}
        =
        \begin{bNiceMatrix}
            \bm{e}_1 & \bm{e}_2 & \cdots & \bm{e}_n
        \end{bNiceMatrix}
        \bm{P}
    \]
    其中$\bm{P}$称为由基$\bm{e}_1, \bm{e}_2, \cdots, \bm{e}_n$到基$\bm{e}'_1, \bm{e}'_2, \cdots, \bm{e}'_n$的\textbf{过渡矩阵}．\\
    向量$\bm{\alpha} \in \bm{V}$，则向量$\bm{\alpha}$在两组基下的坐标满足关系
    \[
        \begin{bNiceMatrix}
            k'_1 \\ k'_2 \\ \vdots \\ k'_n
        \end{bNiceMatrix}
        =
        \bm{P}^{-1}
        \begin{bNiceMatrix}
            k_1 \\ k_2 \\ \vdots \\ k_n
        \end{bNiceMatrix}
    \]
\end{theorem}
\begin{proof}
    设
    \[
        \bm{P} =
        \begin{bNiceMatrix}
            p_{11} & p_{12} & \cdots & p_{1n} \\
            p_{21} & p_{22} & \cdots & p_{2n} \\
            \vdots & \vdots &        & \vdots \\
            p_{n1} & p_{n2} & \cdots & p_{nn}
        \end{bNiceMatrix}
    \]
    则有
    \[
        \bm{e}'_j = p_{1j}\bm{e}_1 + p_{2j}\bm{e}_2 + \cdots + p_{nj}\bm{e}_n, \quad j=1,2,\cdots,n
    \]
    向量$\bm{\alpha}$在两组基下的坐标分别为
    \[
        \bm{\alpha} = k_1\bm{e}_1 + k_2\bm{e}_2 + \cdots + k_n\bm{e}_n
    \]
    和
    \[
        \bm{\alpha} = k'_1\bm{e}'_1 + k'_2\bm{e}'_2 + \cdots + k'_n\bm{e}'_n
    \]
    将$\bm{e}'_j$的表达式代入上式，得
    \[
        \bm{\alpha} = k'_1(p_{11}\bm{e}_1 + p_{21}\bm{e}_2 + \cdots + p_{n1}\bm{e}_n) + k'_2(p_{12}\bm{e}_1 + p_{22}\bm{e}_2 + \cdots + p_{n2}\bm{e}_n) + \cdots + k'_n(p_{1n}\bm{e}_1 + p_{2n}\bm{e}_2 + \cdots + p_{nn}\bm{e}_n)
    \]
    即
    \[
        \bm{\alpha} = (p_{11}k'_1 + p_{12}k'_2 + \cdots + p_{1n}k'_n)\bm{e}_1 + (p_{21}k'_1 + p_{22}k'_2 + \cdots + p_{2n}k'_n)\bm{e}_2 + \cdots + (p_{n1}k'_1 + p_{n2}k'_2 + \cdots + p_{nn}k'_n)\bm{e}_n
    \]
    由向量$\bm{\alpha}$在基$\bm{e}_1, \bm{e}_2, \cdots, \bm{e}_n$下的唯一表示性，得
    \[
        \begin{cases}
            k_1 = p_{11}k'_1 + p_{12}k'_2 + \cdots + p_{1n}k'_n \\
            k_2 = p_{21}k'_1 + p_{22}k'_2 + \cdots + p_{2n}k'_n \\
            \vdots \\
            k_n = p_{n1}k'_1 + p_{n2}k'_2 + \cdots + p_{nn}k'_n
        \end{cases}
    \]
    即
    \[
        \begin{bNiceMatrix}
            k_1 \\ k_2 \\ \vdots \\ k_n
        \end{bNiceMatrix}
        =
        \bm{P}
        \begin{bNiceMatrix}
            k'_1 \\ k'_2 \\ \vdots \\ k'_n
        \end{bNiceMatrix}
    \]
    两边乘以$\bm{P}^{-1}$，得
    \[
        \begin{bNiceMatrix}
            k'_1 \\ k'_2 \\ \vdots \\ k'_n
        \end{bNiceMatrix}
        =
        \bm{P}^{-1}
        \begin{bNiceMatrix}
            k_1 \\ k_2 \\ \vdots \\ k_n
        \end{bNiceMatrix}
    \]
\end{proof}
\begin{example}
    下列说法错误的是（　　）．
    \begin{tasks}[style = itemize,label = \Alph*.](1)
        \task $r<s$时，向量组$\bm{\alpha}_1, \bm{\alpha}_2, \cdots, \bm{\alpha}_r$不能被向量组$\bm{\beta}_1, \bm{\beta}_2, \cdots, \bm{\beta}_s$线性表示
        \task $\bm{\beta}$可由向量组$\bm{\alpha}_1, \bm{\alpha}_2, \cdots, \bm{\alpha}_n$线性表示，则$\bm{\beta}=k_1\bm{\alpha}_1 + k_2\bm{\alpha}_2 + \cdots + k_r\bm{\alpha}_r$，且$k_1,k_2,\cdots,k_r$不全为零
        \task 矩阵$\bm{A}$与矩阵$\bm{B}$等价，则$\bm{A}$的列向量组与$\bm{B}$的列向量组等价
        \task 向量组$\bm{\alpha}_1, \bm{\alpha}_2, \cdots, \bm{\alpha}_n$两两正交，则其为$\R^n$的一组基
    \end{tasks}
\end{example}
\begin{solution}
    答案：ABCD\\
    B选项：取$\bm{\beta} = \bm{0}$\\
    C选项：矩阵等价只能说明其秩相等，不能说明列空间相等．\\
    D选项：零向量与任意向量正交，但不能构成基．
\end{solution}
\section{四大基本空间}
\begin{definition}
    设$\bm{A}$为$m\times n$矩阵，则由$\bm{A}$的列向量所生成的子空间，称为$\bm{A}$的列空间，记为$C(\bm{A})$．
\end{definition}
\begin{definition}
    设$\bm{A}$为$m\times n$矩阵，则由$\bm{A}$的行向量所生成的子空间，称为$\bm{A}$的行空间，记为$R(\bm{A})$或$C(\bm{A}^T)$．
\end{definition}
\begin{definition}
    设$\bm{A}$为$m\times n$矩阵，则线性方程组$\bm{A}\bm{x}=\bm{0}$的解空间，称为$\bm{A}$的零空间，记为$N(\bm{A})$．
\end{definition}
\begin{definition}
    设$\bm{A}$为$m\times n$矩阵，则线性方程组$\bm{A}^T\bm{y}=\bm{0}$的解空间，称为$\bm{A}$的左零空间，记为$N(\bm{A}^T)$．
\end{definition}
\begin{theorem}[秩-零化度定理]
    设$\bm{A}$为$m\times n$矩阵，则有
    \[
        \rank(\bm{A}) + \operatorname{dim}N(\bm{A}) = n
    \]
    \[
        \rank(\bm{A}) + \operatorname{dim}N(\bm{A}^T) = m
    \]
\end{theorem}

\ifx\allfiles\undefined
\end{document}
\fi